\documentclass[10pt,a4paper]{amsart}
\usepackage[margin=19mm]{geometry}
\usepackage{amsmath,amssymb,mathtools}
\usepackage[scr=rsfs,cal=euler]{mathalfa}
\usepackage{xfrac}
\usepackage[unicode,hidelinks,bookmarks=false]{hyperref}
\newcommand{\ie}{i.e.\ }
\newcommand{\cf}{cf.\ }
\newcommand{\resp}{respectively\ }
\newcommand{\Subsection}{\subsection}
\providecommand{\nobreakdash}{\nobreak\hbox{-}}
\setlength{\emergencystretch}{2em}
\multlinegap0pt
\usepackage{bm}
\usepackage{bbm}
\usepackage{dsfont}
\usepackage{xspace}
\usepackage{cancel}
\usepackage{mathtools}
\usepackage{amsthm}
\makeatletter
\@ifundefined{prop}{\newtheorem{prop}{Proposition}}{}
\@ifundefined{lemma}{\newtheorem{lemma}[prop]{Lemma}}{}
\makeatother
\def\BZ{\mathbb{Z}}
\def\CE{\mathcal{E}}
\newcommand{\ord}{\mathrm{ord}}
\title[Divisibility of the coefficients of modular polynomials]{Divisibility of the coefficients of modular polynomials}
\author{Florian Breuer}
\date{Source: arXiv:2509.06423v3 (CC BY 4.0)}
\begin{document}
\maketitle

\noindent\emph{Source and licence.} Divisibility of the coefficients of modular polynomials. Version arXiv:2509.06423v3 (CC BY 4.0), distributed under the Creative Commons Attribution 4.0 International licence, as recorded in the arXiv licence metadata for this exact version and shown on the abstract page https://arxiv.org/abs/2509.06423v3. Full rights notice: licenses/arxiv-2509.06423-CC-BY-4.0.txt.\par\smallskip

\noindent\textit{Editorial scope.} Counted unit: the Prop whose source label is \texttt{prop:general} in the article (source0.tex lines 530--547, source label \texttt{prop:general}), delivered together with the article's own complete original proof of it (source0.tex lines 549--569, one \texttt{proof} environment of 21 lines). No mathematics is written, repaired or re-derived: statement and proof are the article's own text line for line. Required context retained uncounted: lem:interpolation (lines 496--509). Every definition, display and label of those lines is retained unchanged. Omitted from this package and not redistributed: 1--495, 510--529, 548--548, 570--724. Nothing inside a retained excerpt is omitted or altered: every retained body is delivered exactly as the article's lines are written, so no omission receipt of any kind accompanies this package. Citations: the retained excerpts carry no citation command at all, so no bibliography entry of the article is transcribed and no reference is redistributed. Reference adaptation: none was needed and none was made; every reference inside the delivered unit resolves inside it, so no unresolved reference or citation is produced. Printed slips kept literal: no printed word, symbol, hypothesis or number of the counted statement or of its proof was altered. Layout and class: the article's own class is \texttt{amsart} with packages including fontenc, amsmath, amssymb, amsfonts, amsthm, hyperref; the wrapper is the lane's pinned \texttt{amsart} editorial wrapper at 10pt on A4, which restates the theorem-like environments the retained text needs and adds the article's own macro definitions that the retained text needs (\texttt{\textbackslash BZ}, \texttt{\textbackslash CE}, \texttt{\textbackslash ord}), with extra wrapper packages (bm, bbm, dsfont, xspace, cancel, mathtools). The delivered numbering is the unit's own. Assets: no retained span contains a figure environment, a graphics-inclusion command, a PDF-page inclusion, a table or a TikZ picture; no image file is copied into this package, the package declares no asset dependency and no image file is redistributed with it. Licence: version arXiv:2509.06423v3 of this article is distributed under the Creative Commons Attribution 4.0 International licence, as recorded in the arXiv licence metadata for this exact version and shown on the abstract page https://arxiv.org/abs/2509.06423v3. A full rights notice is delivered at licenses/arxiv-2509.06423-CC-BY-4.0.txt. Characters: every retained excerpt is pure ASCII, so the delivered engine, XeLaTeX with its default Latin Modern text font, reports no missing character.
\begin{lemma}\label{lem:interpolation}
    Let $f(Y) = a_0 + a_1Y + \ldots + a_dY^d \in K[Y]$.
    Fix $n\in\BZ$ and
    let $y_0, y_1,\ldots, y_d\in K$ be such that
    \begin{enumerate}
        \item $v(y_0)=v(y_1)=\cdots=v(y_d)=n$,
        \item $v(y_k-y_l) = n$ for all $k\neq l$.
    \end{enumerate}
    Then 
    \[
    v(a_j)  \geq \min_{0\leq k \leq d} v(f(y_k)) - nj
    \quad \text{for all $j=0,1,2,\ldots, d$}.
    \]
\end{lemma}

\begin{prop}\label{prop:general}
	Let $E_1$ and $E_2$ be elliptic curves over $K$ with good reduction and $j$-invariants $J_1 = j(E_1)$ and $J_2 = j(E_2)$.
    Let $N > 1$ with $v(N)=0$.
    Fix positive integers $n_1$ and $n_2$, and suppose that there exist elliptic curves $\CE_k/K$, $k=0,1,\ldots, \psi(N)$ satisfying the following conditions:
    \begin{enumerate}
        \item Each $\CE_k$ has good reduction;
        \item $v(j(\CE_k)-J_2) = v(j(\CE_k)-j(\CE_l)) = n_2$ for all $k\neq l$;
        \item For every $k$ and every elliptic curve $\tilde{\CE}_k$ linked to $\CE_k$ by a cyclic isogeny of degree $N$, we have
        \[
        v(j(\tilde{\CE}_k) - J_1) > 0 \Longrightarrow v(j(\tilde{\CE}_k) - J_1) \geq n_1.
        \]
    \end{enumerate}
    Then the coefficients of $\Phi_N(X+J_1, Y+J_2) = \sum_{0\leq i,j \leq \psi(N)}a_{i,j}X^iY^j\in A[X,Y]$ satisfy
    \[
    v(a_{i,j}) \geq n_1 C_{J_1,J_2}(N,\pi)-n_1 i-n_2 j,
    \]
    where $C_{J_1,J_2}(N,\pi) = \ord_X\big(\Phi_N(X+J_1,J_2) \bmod \pi\big).$
\end{prop}

\begin{proof}
	For ease of notation, let $C = C_{J_1,J_2}(N,\pi).$
Write 
\begin{align*}
    & \Phi_N(X+J_1,Y+J_2) = b_0(Y) + b_1(Y)X + \cdots + b_{\psi(N)-1}(Y)X^{\psi(N)-1} + X^{\psi(N)} \\
    & b_i(Y) = a_{i,0} + a_{i,1}Y + \cdots + a_{i,\psi(N)}Y^{\psi(N)}, \quad i=0,\ldots, \psi(N).
\end{align*}

For each $k$, let $y_k = j(\CE_k) - J_2$. The roots of $\Phi_N(X, y_k + J_2)$ are the $j$-invariants of elliptic curves $\CE_{k,m}$ linked to $\CE_k$ by a cyclic $N$-isogeny. Since $v(N)=0$, $\CE_k[N]$ is unramified and these elliptic curves and isogenies are all defined over $K$, since unramified extensions of $K$ correspond to extensions of the residue field $A/\pi$, which is algebraically closed.

By definition, $C$ of these roots $j(\CE_{k,m})$, $m=1,2,\ldots, C$, satisfy 
$v(j(\CE_{k,m})-J_1) >0$, so by assumption $v(j(\CE_{k,m} )-J_1) \geq n_1$, and the rest satisfy $v(j(\CE_{k,m})-J_1)=0$. 

The coefficients $b_i(y_k)$ of 
$\Phi_N(X+J_1, y_k+J_2)$ are symmetric forms in the roots $j(\CE_{k,m})-J_1$ which satisfy $v(j(\CE_{k,m})-J_1) \geq n_1$ for $m=1,2,\ldots, C$, so it follows that
\[
v(b_i(y_k)) \geq n_1(C -i), \quad i = 0, 1, \ldots, C.
\]

Now applying Lemma \ref{lem:interpolation} to the polynomials $b_i(Y)$ with the interpolation points $y_k$ completes the proof.
\end{proof}

\end{document}
